%! AP Statistics — Unit 2, Lesson 2.4 — Slides
\documentclass[aspectratio=169, 11pt]{beamer}
\usepackage{apstats-beamer}

\begin{document}

% ── TITLE SLIDE ───────────────────────────────────────────────────────────────
{
\setbeamercolor{background canvas}{bg=navy}
\begin{frame}[plain]
  \vspace{2.8cm}\hspace{0.55cm}%
  \begin{minipage}{0.9\textwidth}
    {\color{goldacc}\bfseries\small LESSON 2.4}\\[0.5cm]
    {\color{white}\bfseries\LARGE Representing the Relationship Between Two Quantitative Variables}\\[1.2cm]
    {\color{skymid}\small AP Statistics~$\cdot$~Unit 2: Exploring Two-Variable Data}
  \end{minipage}
\end{frame}
}

% ── LEARNING TARGETS ─────────────────────────────────────────────────────────
\begin{frame}[t, plain]
  \navyheader{Today's Learning Targets}
  \vspace{0.3cm}
  By the end of class, I will be able to\dots
  \begin{itemize}
    \item identify the \textbf{explanatory} and \textbf{response} variables and place them on the correct axes.
    \item \textbf{construct} a correctly labeled scatterplot from a two-variable dataset.
    \item describe the \textbf{direction}, \textbf{form}, \textbf{strength}, and \textbf{outliers} of a scatterplot in context.
    \item write a complete \textbf{DOFS description} using the DOFS framework.
  \end{itemize}
  \vspace{0.4cm}
  \textcolor{goldacc}{\small Standards: UNC-1.S [Skill 2.A] $\cdot$ UNC-1.T [Skill 4.B]}
\end{frame}

% ── HOOK ─────────────────────────────────────────────────────────────────────
\begin{frame}[t, plain]
  \navyheader{Hook --- Which Pattern Is Most Useful?}
  \vspace{0.3cm}
  \textit{Four scatterplots are displayed. Rank them from most useful for prediction to least useful.}
  \vspace{0.4cm}
  \begin{columns}
    \column{0.48\textwidth}
      \textbf{What makes a pattern useful?}
      \begin{itemize}
        \item Points follow a clear, consistent shape.
        \item We can predict $y$ reliably from $x$.
        \item The scatter around the pattern is small.
      \end{itemize}
    \column{0.48\textwidth}
      \textbf{What makes a pattern less useful?}
      \begin{itemize}
        \item Points are scattered randomly.
        \item The relationship is curved in an irregular way.
        \item A few extreme points distort the picture.
      \end{itemize}
  \end{columns}
  \vspace{0.5cm}
  \textcolor{navy}{\small Today we learn the precise language to describe these patterns: \textbf{DOFS}.}
\end{frame}

% ── WHAT IS A SCATTERPLOT ────────────────────────────────────────────────────
\begin{frame}[t, plain]
  \navyheader{What Is a Scatterplot?}
  \vspace{0.3cm}
  A \textbf{scatterplot} displays two \textbf{quantitative} variables for the same set of individuals.
  \vspace{0.4cm}
  \begin{columns}
    \column{0.48\textwidth}
      \textbf{Axis rule:}
      \begin{itemize}
        \item \textbf{Explanatory variable} ($x$) on the horizontal axis.
        \item \textbf{Response variable} ($y$) on the vertical axis.
        \item Each point = one individual.
      \end{itemize}
    \column{0.48\textwidth}
      \textbf{Label your axes:}
      \begin{itemize}
        \item Variable name \textit{and} units.
        \item Consistent, reasonable scale.
        \item Start at a value that shows the data well.
      \end{itemize}
  \end{columns}
\end{frame}

% ── DOFS: DIRECTION ──────────────────────────────────────────────────────────
\begin{frame}[t, plain]
  \navyheader{DOFS: \textbf{D}irection}
  \vspace{0.3cm}
  \begin{columns}
    \column{0.32\textwidth}
      \textbf{Positive}\\
      As $x$ increases, $y$ tends to \textbf{increase}.\\[0.2cm]
      \textit{Points go up from left to right.}
    \column{0.32\textwidth}
      \textbf{Negative}\\
      As $x$ increases, $y$ tends to \textbf{decrease}.\\[0.2cm]
      \textit{Points go down from left to right.}
    \column{0.32\textwidth}
      \textbf{None}\\
      Knowing $x$ gives \textbf{no information} about $y$.\\[0.2cm]
      \textit{Points show no slope.}
  \end{columns}
  \vspace{0.8cm}
  \textcolor{navy!70}{\small Caution: ``positive'' does not mean ``good'' --- it only describes the direction of the pattern.}
\end{frame}

% ── DOFS: FORM AND STRENGTH ──────────────────────────────────────────────────
\begin{frame}[t, plain]
  \navyheader{DOFS: \textbf{F}orm and \textbf{S}trength}
  \vspace{0.3cm}
  \begin{columns}
    \column{0.48\textwidth}
      \textbf{Form} --- overall shape:
      \begin{itemize}
        \item \textbf{Linear} --- points cluster around a straight line.
        \item \textbf{Curved} --- points follow a curved path.
        \item \textbf{None} --- no discernible pattern.
      \end{itemize}
      \vspace{0.3cm}
      \textit{Form matters: correlation and regression only apply to \textbf{linear} associations.}
    \column{0.48\textwidth}
      \textbf{Strength} --- how closely points follow the form:
      \begin{itemize}
        \item \textbf{Strong} --- points hug the form tightly.
        \item \textbf{Moderate} --- points follow the form with some scatter.
        \item \textbf{Weak} --- points are loosely scattered.
      \end{itemize}
  \end{columns}
\end{frame}

% ── DOFS: OUTLIERS ───────────────────────────────────────────────────────────
\begin{frame}[t, plain]
  \navyheader{DOFS: \textbf{O}utliers}
  \vspace{0.3cm}
  An \textbf{outlier} in a scatterplot is a point that \textbf{deviates from the overall pattern}.
  \vspace{0.4cm}
  \begin{columns}
    \column{0.48\textwidth}
      A bivariate outlier may have:
      \begin{itemize}
        \item An extreme $x$ value.
        \item An extreme $y$ value.
        \item Both $x$ and $y$ in range but \textit{far from the pattern}.
      \end{itemize}
    \column{0.48\textwidth}
      \textbf{Describe it in context:}\\[0.2cm]
      ``One [individual] with [x-value] had [y-value], which does not follow the overall [direction] pattern.''
  \end{columns}
  \vspace{0.5cm}
  \textcolor{navy!70}{\small Preview: Lesson 2.8 --- \textit{influential points} are outliers that substantially change the regression line when removed.}
\end{frame}

% ── PUTTING IT TOGETHER ───────────────────────────────────────────────────────
\begin{frame}[t, plain]
  \navyheader{Putting It Together: The DOFS Template}
  \vspace{0.3cm}
  \textbf{Complete description template:}\\[0.4cm]
  \begin{block}{}
    ``There is a \textbf{[strength]} \textbf{[direction]} \textbf{[form]} association between
    \textbf{[explanatory variable]} and \textbf{[response variable]} for these \textbf{[individuals]}.
    \textbf{[Outlier sentence if applicable.]}''
  \end{block}
  \vspace{0.4cm}
  \textbf{Example (sleep vs.\ reaction time):}\\[0.2cm]
  ``There is a strong negative linear association between hours of sleep and reaction time (ms)
  for these 10 students. As sleep increases, reaction time tends to decrease. There are no
  obvious outliers.''
\end{frame}

% ── CLOSING ──────────────────────────────────────────────────────────────────
{
\setbeamercolor{background canvas}{bg=navy}
\begin{frame}[plain]
  \vspace{3cm}\hspace{0.55cm}%
  \begin{minipage}{0.9\textwidth}
    {\color{goldacc}\bfseries\large Exit Ticket}\\[0.6cm]
    {\color{white} A scatterplot shows hours of TV per day vs.\ GPA for 25 students.
    The points slope downward from left to right, with most clustering tightly along the line.
    One student who watches 7 hr/day has a GPA of 3.8.}\\[0.5cm]
    {\color{skymid}\small
      (1) Identify the explanatory and response variables.\\
      (2) Write a complete DOFS description in context.\\
      (3) Describe the red point and explain whether it is an outlier.}
  \end{minipage}
\end{frame}
}

\end{document}
